\chapter*{چکیدهٔ اجرایی}
\addcontentsline{toc}{chapter}{چکیدهٔ اجرایی}

پیشرفت روزافزون رایانش کوانتومی، به‌ویژه بلوغ الگوریتم‌های تجزیهٔ اعداد صحیح نظیر الگوریتم \lr{\cite{shor1997}}، امنیت زیرساخت‌های رمزنگاری کلید عمومی کنونی -- شامل \lr{RSA}\LTRfootnote{Rivest--Shamir--Adleman} و \lr{ECC}\LTRfootnote{Elliptic Curve Cryptography} -- را که ستون فقرات تراکنش‌های بانکی، امضای دیجیتال و ارتباطات محرمانهٔ دولتی هستند، در افق یک تا دو دهه با تهدید جدی مواجه می‌کند. این تهدید تنها به آینده محدود نمی‌شود؛ حملات موسوم به «برداشت اکنون، رمزگشایی بعداً» \lr{(Harvest Now, Decrypt Later)} به این معناست که داده‌های رمزشدهٔ امروز، در صورت رهگیری و ذخیره‌سازی، می‌توانند در آینده توسط یک رایانهٔ کوانتومی رمزگشایی شوند. این پروپوزال دو راهکار مکمل و در حال بلوغ صنعتی را برای مقابله با این تهدید بررسی می‌کند: نخست، توزیع کلید کوانتومی \lr{(QKD)} که با تکیه بر قوانین بنیادین فیزیک کوانتومی -- به‌ویژه قضیهٔ عدم-شبیه‌سازی \lr{\cite{wootters1982}} -- امنیت تبادل کلید را نه بر پیچیدگی محاسباتی، بلکه بر قوانین فیزیک استوار می‌کند؛ و دوم، رمزنگاری پساکوانتومی \lr{(PQC)} که با استانداردهای تازه نهایی‌شدهٔ \lr{NIST}\LTRfootnote{National Institute of Standards and Technology (United States)} یعنی \lr{ML-KEM (FIPS 203)}، \lr{ML-DSA (FIPS 204)} و \lr{SLH-DSA (FIPS 205)} \lr{\cite{nist2024release}} امکان مهاجرت نرم‌افزاری و مقیاس‌پذیر را فراهم می‌آورد.

در این سند، پس از مرور مبانی نظری و معماری‌های فنی \lr{QKD} (فصل‌های \lr{2} و \lr{3}) و مرور استانداردهای \lr{PQC} (فصل \lr{4})، رویکردی \textbf{ترکیبی و لایه‌ای} \lr{(Hybrid QKD+PQC)} برای تأمین امنیت دفاع در عمق پیشنهاد می‌شود. کاربردهای واقعی این فناوری‌ها در بخش بانکی -- از جمله تجربهٔ \lr{HSBC} در تأمین امنیت معاملات ارزی و طلای توکنیزه‌شده، و شبکهٔ \lr{Q-CAN} بانک \lr{JPMorgan Chase} -- به‌همراه تجربهٔ کشورهایی نظیر چین (کریدور کوانتومی پکن--شانگهای)، سوئد (پروژهٔ \lr{EuroQCI})، سنگاپور (\lr{NQSN+}) و کرهٔ جنوبی مرور می‌شود (فصل‌های \lr{6} و \lr{7}). سپس یک طرح فنی مشخص برای ایجاد یک حلقهٔ کلان‌شهری \lr{QKD} در تهران -- با گره‌های پیشنهادی در بانک مرکزی، مراکز دادهٔ بانک‌های بزرگ، شرکت مخابرات ایران و مراکز دولتی حساس -- همراه با فهرست تجهیزات، برآورد هزینه (فصل \lr{9}) و نقشهٔ راه اجرایی سه‌مرحله‌ای در افق \lr{18} تا \lr{36} ماه ارائه می‌گردد. تحلیل ریسک امنیتی (فصل \lr{10}) به‌طور ویژه به آسیب‌پذیری معماری گره مورد اعتماد \lr{(Trusted Node)} و راهکارهای کاهش آن از طریق پروتکل‌های \lr{MDI-QKD}\LTRfootnote{Measurement-Device-Independent QKD} و لایهٔ ترکیبی \lr{PQC} می‌پردازد.

هدف نهایی این طرح، ایجاد یک زیرساخت پایلوت ملی است که علاوه‌بر ارتقای تاب‌آوری سایبری نظام بانکی و زیرساخت‌های حیاتی کشور در برابر تهدیدات کوانتومی، بستری برای توسعهٔ بومی دانش فنی، تربیت نیروی متخصص و کاهش وابستگی به فناوری‌های خارجی فراهم آورد.

